\color{unimain}{\subsection{Дифференциальная защита ошиновки}}
\color{black}

\begin{enumerate}

%%%%%%%%%%%%%%%%%%%%%%%%%%%%%%%%%%%%%%%%%%%%%%%%%% ОБШИЕ СВЕДЕНИЯ %%%%%%%%%%%%%%%%%%%%%%%%%%%%%%%%%%%%%%%%%%%%%%%%%%
\item \textbf{Общие сведения}

\begin{enumerate}

\item
ДЗО является быстродействующей защитой с абсолютной селективностью и выполняет функцию основной токовой защиты ошиновки от всех видов КЗ.

\item
ДЗО в составе данного устройства имеет трехфазное исполнение, т.е. терминал обрабатывает все три фазы с каждого присоединения. Ввод и вывод данной функции осуществляется либо посредством одного программного ключа \textbf{«ДЗО»}, либо с использованием внешнего сигнала \textit{«ДЗО: Вывод»}.

\item
Параметры для настройки ДЗО приведены в таблице \ref{app_set:dzo}.

\item
В устройстве предусмотрена настройка полярности ТТ присоединений. По умолчанию измерение токов выполняется в направлении к защищаемому объекту, то есть все ТТ считаются включёнными встречно друг к другу. Если для какого-либо присоединения фактическая полярность не совпадает с принятой по умолчанию, пользователь может скорректировать её двумя способами: либо изменением подключения вторичных токовых цепей, либо установкой программного ключа \textbf{«Инверсия Bn»} в положение \textbf{«Введено»}.

\end{enumerate}

%%%%%%%%%%%%%%%%%%%%%%%%%%%%%%%%%%%%%%%%%%%%%%%%%% Измерительные органы ДЗО %%%%%%%%%%%%%%%%%%%%%%%%%%%%%%%%%%%%%%%%%%%%%%%%%%
\item \textbf{Измерительные органы ДЗО}

\begin{enumerate}

\item
Дифференциальная защита состоит из следующих измерительных органов:
\begin{itemize} 
\item дифференциальная токовая отсечка (ДТО);
\item дифференциальный орган с торможением (ДОТХ);
\item блокировка при внешних КЗ (БВКЗ);
\item чувствительный токовый орган (ЧТО);
\item орган блокировки по второй гармонике дифференциального тока;
\item орган дифференциального тока для КЦТ.
\end{itemize}

\item
Алгоритм работы измерительной части ДЗО фазы А приведен на рисунке \ref{pict_dzo:IO_DZO}, алгоритмы работы измерительной части других фаз выполнены аналогично.

\medskip
\begin{figure}[H]
\centering
\includegraphics[width=0.8\textwidth,height=0.8\textheight,keepaspectratio]{\fbpath/05.01 ДЗО/661.0641 ДЗО/img/IO_DZO.pdf}
\captionof{figure}{Функциональная схема логики измерительной части ДЗО ф.А}\label{pict_dzo:IO_DZO}
\end{figure}

\end{enumerate}

%%%%%%%%%%%%%%%%%%%%%%%%%%%%%%%%%%%%%%%%%%%%%%%%%% Дифференциальная токовая отсечка %%%%%%%%%%%%%%%%%%%%%%%%%%%%%%%%%%%%%%%%%%%%%%%%%%
\item \textbf{Дифференциальная токовая отсечка}

\begin{enumerate}

\item
ДТО предназначена для быстрого отключения защищаемого объекта при внутренних повреждениях, сопровождающихся большими токами. Ввод в работу данной функции для всех дифференциальных органов осуществляется с помощью одного программного ключа \textbf{«ДТО»}.

\item
Данный ИО выполнен без торможения и может быть заблокирован только при обнаружении неисправности в токовых цепях. В качестве рабочей величины использует действующее значение первой гармоники дифференциального тока, которое определяется как модуль геометрической суммы первых гармоник токов присоединений, входящих в зону действия защиты:
\begin{align}
I_{a\textmd{д}}=\mid \sum_{n=1}^N \dot{I}_{a,\,n}^{\,'} \mid, \\
I_{b\textmd{д}}=\mid \sum_{n=1}^N \dot{I}_{b,\,n}^{\,'} \mid, \\
I_{c\textmd{д}}=\mid \sum_{n=1}^N \dot{I}_{c,\,n}^{\,'} \mid, 
\end{align}
\begin{eqexpl} [15mm]
    \item{$ I_{a\textmd{д}}, I_{b\textmd{д}}, I_{c\textmd{д}} $} действующие значения первых гармоник дифференциальных токов фаз A, B и C, о.е.;
	\item{$ N $} количество присоединений;
	\item{$ \dot{I}_{a,\,n}^{\,'}, \dot{I}_{b,\,n}^{\,'}, \dot{I}_{c,\,n}^{\,'} $} вектора первых гармоник выровненных фазных токов n-го присоединения, о.е.
\end{eqexpl}

\item
Ток срабатывания ДТО задается уставкой \textbf{«Iдто»}.

\end{enumerate}


%%%%%%%%%%%%%%%%%%%%%%%%%%%%%%%%%%%%%%%%%%%%%%%%%% Дифференциальный орган с торможением %%%%%%%%%%%%%%%%%%%%%%%%%%%%%%%%%%%%%%%%%%%%%%%%%%
\item \textbf{Дифференциальный орган с торможением}

\begin{enumerate}

\item
ДОТХ является основным измерительным органом и предназначен для отключения защищаемого объекта при всех видах внутренних повреждений.

\item
В качестве рабочих величин данный ИО использует действующие значения дифференциального и тормозного токов первой гармоники. Действующее значение тормозного тока первой гармоники рассчитывается как полусумма модулей токов присоединений, входящих в зону действия защиты:
\begin{align}
I_{a\textmd{т}}=0,5 \cdot \sum_{n=1}^N \mid \dot{I}_{a,\,n}^{\,'} \mid, \\
I_{b\textmd{т}}=0,5 \cdot \sum_{n=1}^N \mid \dot{I}_{b,\,n}^{\,'} \mid, \\
I_{c\textmd{т}}=0,5 \cdot \sum_{n=1}^N \mid \dot{I}_{c,\,n}^{\,'} \mid
\end{align}

\item
Тормозная характеристика дифференциальной защиты имеет два участка. Горизонтальный участок не зависит от тормозного тока и необходим для отстройки от небольших токов небаланса при протекании малых сквозных токов, которые не вызывают насыщение ТТ. Уровень срабатывания данного задается отдельной уставкой: \textbf{«Iд нач.»}. Наклонный участок обеспечивает отстройку от токов небаланса при насыщении ТТ.
Наклон характеристики определяется коэффициентом торможения, а точка начала наклона – значением начального тормозного тока. Эти параметры задаются отдельными уставками: \textbf{«Kт»} и \textbf{«Iт нач.»}.

\medskip
\begin{figure}[H]
\centering
\includegraphics[width=0.55\textwidth,height=0.55\textheight,keepaspectratio]{\fbpath/05.01 ДЗО/661.0641 ДЗО/img/DiffZ_torm_char.pdf}
\captionof{figure}{Тормозная характеристика дифференциальной защиты}\label{pic_dzo:DiffZ_torm_char}
\end{figure}
\medskip

\end{enumerate}


%%%%%%%%%%%%%%%%%%%%%%%%%%%%%%%%%%%%%%%%%%%%%%%%%% Орган блокировки при внешнем КЗ %%%%%%%%%%%%%%%%%%%%%%%%%%%%%%%%%%%%%%%%%%%%%%%%%%
\item \textbf{Орган блокировки при внешнем КЗ}

\begin{enumerate}

\item
БВКЗ предназначена для выявления внешнего КЗ и обеспечивает время достоверного измерения тока на уровне 5 мс. Ввод в работу данной функции осуществляется с помощью программного ключа \textbf{«БВКЗ»}.

\item
Принцип действия органа основан на сравнении направлений токов присоединений. Для идентификации внешнего повреждения формируются два вектора: вектор тока повреждённого присоединения ($ \dot{I}_1 $), вектор тока, рассчитываемого как разность между дифференциальным током и током повреждения ($ \dot{I}_2 $), т.е. вектор суммы токов всех остальных присоединений.

\item
В качестве вектора тока поврежденного присоединения выступает комплексное действующее значение основной гармоники тока того присоединения, в котором зафиксировано максимальное значение тока ($ \dot{I} \textit{макс} $).

\begin{align}
\begin{split}
\dot{I}_1&=\dot{I}_\textit{макс}\\
\dot{I}_2&=\sum_{n=1}^N \dot{I}_{n} - \dot{I}_\textit{макс}
\end{split}
\end{align}

\item
Выбор присоединения с максимальным током основана на сравнении мгновенных значений токов всех присоединений. Если наибольший из этих токов превышает уставку \textbf{«I БВКЗ»}, фиксируется номер соответствующего присоединения и его комплексное действующее значение основной гармоники тока. Указанное присоединение принимается в качестве поврежденного и сохраняется до тех пор, пока действующее значение его тока превышает указанную выше уставку.

\item
БВКЗ анализирует угол между двумя сформированными векторами. При внешнем КЗ угол между векторами превышает 90º. Данный признак используется как критерий для формирования сигнала блокировки.

\medskip
\begin{figure}[H]
\centering
\begin{subfigure}{0.4\textwidth}
    \centering
    \includegraphics[width=\textwidth, trim=0 0 0 0, clip]{\fbpath/05.01 ДЗО/661.0641 ДЗО/img/BVKZ_a.pdf}
    \caption{\textit{а}}
    \label{pict_dzo:Action_DFZ_a}
\end{subfigure}
%\hfill
\begin{subfigure}{0.4\textwidth}
    \centering
    \includegraphics[width=\textwidth, trim=0 0 0 0, clip]{\fbpath/05.01 ДЗО/661.0641 ДЗО/img/BVKZ_b.pdf}
    \caption{\textit{б}}
    \label{pict_dzo:Action_DFZ_b}
\end{subfigure}

\caption{Характеристика срабатывания БВКЗ: \\ \textit{а} -- внутреннее КЗ; \textit{б} -- внешнее КЗ}
\label{pict_dzo:BVKZ}
\end{figure}
\medskip

\end{enumerate}


%%%%%%%%%%%%%%%%%%%%%%%%%%%%%%%%%%%%%%%%%%%%%%%%%% Чувствительный токовый орган %%%%%%%%%%%%%%%%%%%%%%%%%%%%%%%%%%%%%%%%%%%%%%%%%%
\item \textbf{Чувствительный токовый орган}

\begin{enumerate}

\item
В устройстве предусмотрено два подхода для надежного отключения выключателей присоединений при работе ДЗО, в том числе в режимах АПВ шин и опробования. Такие режимы могут сопровождаться пониженным уровнем токов КЗ, в результате чего чувствительность дифференциальных органов может быть недостаточной. Первый подход подразумевает применение ЧТО. Второй подход заключается в очувствлении уставок тормозной характеристики, которые используются в ДОТХ.

\item
ЧТО выполнен в виде реле тока, включенного на дифференциальный ток. Ток срабатывания ЧТО задается отдельной уставкой \textbf{«Iчто»}.

\item
Очувствление представляет собой автоматическую смену величины начального тока срабатывания тормозной характеристики на более чувствительную уставку. Параметры тормозной характеристики в режиме очувствления задаются отдельными уставками. Подробное описание алгоритма очувствления приведено в подразделе \ref{logic_och}.

\end{enumerate}


%%%%%%%%%%%%%%%%%%%%%%%%%%%%%%%%%%%%%%%%%%%%%%%%%% Блокировка по второй гармонике дифференциального тока %%%%%%%%%%%%%%%%%%%%%%%%%%%%%%%%%%%%%%%%%%%%%%%%%%
\item \textbf{Блокировка по второй гармонике дифференциального тока}

\begin{enumerate}

\item
Для предотвращения срабатывания ИО ДЗО (за исключением ДТО) при БНТ, возникающих при включении присоединения с силовым трансформатором, используется блокировка по второй гармонике дифференциального тока. Влияние второй гармоники, возникающей при БНТ, на работу ИО ДЗО проявляется только в условиях опробования по схеме с «открытым плечом».

\item
ИО блокировки выполнен в виде реле тока, реагирующее на отношение дифференциального тока второй гармоники к дифференциальному току первой гармоники. Величина срабатывания для каждого дифференциального органа задается отдельной уставкой \textbf{«I2г/I1г ДЗО»}.

\item
В устройстве предусмотрена возможность ввода перекрестной блокировки по второй гармонике. Перекрестная блокировки вводится программным ключом \textbf{«Перекр. блок.»}.  При её вводе срабатывание блокировки по второй гармонике в одной из фаз приводит к блокированию ИО ДЗО в остальных фазах на время, задаваемое уставкой \textbf{«Тперекр. блок.»}. 


\end{enumerate}

%%%%%%%%%%%%%%%%%%%%%%%%%%%%%%%%%%%%%%%%%%%%%%%%%% Формирование сигнала срабатывания ДЗО  %%%%%%%%%%%%%%%%%%%%%%%%%%%%%%%%%%%%%%%%%%%%%%%%%%
\item \textbf{Формирование сигнала срабатывания ДЗО}

\begin{enumerate}

\item
При срабатывании ИО ДЗО, сигнал на отключение сформируется по истечении соответствующей выдержки времени \textbf{«Тср ДЗО»}.

\item
Алгоритм работы функциональной части ДЗО приведен на рисунке \ref{pict_dzo:DZO}.

\medskip
\begin{figure}[H]
\centering
\includegraphics[width=0.9\textwidth,height=0.9\textheight,keepaspectratio]{\fbpath/05.01 ДЗО/661.0641 ДЗО/img/DZO.pdf}
\captionof{figure}{Функциональная схема логики ДЗО}\label{pict_dzo:DZO}
\end{figure}
\medskip

\end{enumerate}


%%%%%%%%%%%%%%%%%%%%%%%%%%%%%%%%%%%%%%%%%%%%%%%%%% Контроль цепей тока %%%%%%%%%%%%%%%%%%%%%%%%%%%%%%%%%%%%%%%%%%%%%%%%%%
\item \textbf{Контроль цепей тока}

\begin{enumerate}

\item
В устройстве реализован контроль исправности токовых цепей для исключения ложного срабатывания ДЗО вследствие нарушения баланса токов в дифференциальной цепи при обрыве, замыканий и т.д.

\item
Алгоритм работы КЦТ приведен на рисунке \ref{pict_dzo:KCT}.

\medskip
\begin{figure}[H]
\centering
\includegraphics[width=0.9\textwidth,height=0.9\textheight,keepaspectratio]{\fbpath/05.01 ДЗО/661.0641 ДЗО/img/KCT.pdf}
\captionof{figure}{Функциональная схема логики КЦТ}\label{pict_dzo:KCT}
\end{figure}
\medskip

\item
Ввод и вывод контроля исправности токовых цепей осуществляется программным ключом \textbf{«КЦТ»}.

\item
Для выявления неисправности токовых цепей предусмотрен измерительный орган, реагирующий на дифференциальный ток. Ток срабатывания КЦТ задается отдельной уставкой \textbf{«Iкцт»}.

\item
При обнаружении небаланса в токовых цепях измерительный орган с соответствующей выдержкой времени \textbf{«Ткцт»} действует на сигнализацию, а также на блокировку ДЗО. Блокировка ДЗО от КЦТ осуществляется при введенном программном ключе \textbf{«Реж КЦТ»}.

\item
Сброс блокировки выполняется внешним сигналом \textit{«ДЗО/ ДЗО: Возврат КЦТ»}. Оперативный вывод блокировки ДЗО при неисправностях в токовых цепях осуществляется внешним сигналом \textit{«ДЗО/ ДЗО: Блок. ДЗО от КЦТ»}.

\end{enumerate}

\end{enumerate}
